\documentclass[11pt,a4paper]{article}
 
\usepackage{amsmath}
\usepackage{authblk}
\usepackage[portuguese]{babel}
\usepackage{enumitem}
\usepackage{eurosym}\def\texteuro{\euro}
\usepackage{fancyhdr}
\usepackage[T1]{fontenc}
\usepackage[hmargin=2cm,vmargin=3.5cm,bmargin=3cm]{geometry}
\usepackage{listings} % for code snippets
\usepackage{graphicx}
\usepackage[utf8]{inputenc}
\usepackage{lastpage}
\usepackage{listings}

% listing environment for code
\lstnewenvironment{code}[1][]%
  {\noindent\minipage{\linewidth}\medskip 
   \lstset{basicstyle=\ttfamily\footnotesize,frame=single,#1}}
  {\endminipage}

\newcommand{\select}{\mbox{\Large$\sigma$}}
%\pagenumbering{arabic}
\title{\textbf{Projecto de Bases de Dados, Parte 2}}

\author{Bruno Cardoso (72619), Lídia Freitas (78559) e Rodrigo Bernardo (78942)} 
\affil{Instituto Superior Técnico}

\pagestyle{fancy}
\cfoot{Página \thepage\ de \pageref{LastPage}}

\begin{document}
\date {\today}

\maketitle
\thispagestyle{empty} % remove numeracao na primeira pagina
\centerline{\includegraphics[width=0.4\textwidth]{ist-simbolo.jpg}}

\begin{description}[noitemsep]
	\item \centering{Grupo 17}
	\item Turno: Quinta-Feira, 08h00, LAB 14
\end{description}

\newpage

\tableofcontents
\newpage

\section{Introdução}

Este relatório foi elaborado na âmbito da cadeira de Bases de Dados e sumariza o trabalho realizado para a segunda parte do projecto da cadeira.

Nesta disciplina, o projecto envolveu criar um Bloco de Notas, ou seja, uma conjunto de procedimentos e interrogações que melhor descrevem uma forma de armazenamento e consulta de informação. Tal como o bloco de notas tradicional, a implementação deste projecto seguiu as várias categorias para organizar a informações, tal como a implementação de registos, páginas, valores. De forma a limitar o espaço de cada utilizador ao seu bloco de notas, foram criadas diversas verificações para o acesso á informação, garantindo que um utilizador só pode ver informações a que tenha permissão para tal.
	Este projecto foi implementado utilizando o SGBD Mysql, e as diversas páginas de acesso á informação foram criadas utilizando PHP. Foram tomadas as devidas precauções no que toca ao tratamento de inputs, de forma a manter a segurança da base de dados prevenida, ao fazer a retirada de caracteres indesejados no input.
	Todos os acessos á base de dados são atómicos, garantido que cada acesso é unico ao mesmo tempo.



Neste pretendemos mostrar a nossa implementação de um bloco de notas assim como os aspectos de interação e melhoria da mesma.


\newpage

\section{Consultas SQL}
% http://sqlformat.org/
\begin{enumerate}[label=(\alph*)]
\begin{minipage}{\textwidth}
	\item Quais são os utilizadores que falharam o login mais vezes do que tiveram sucesso?
	
		\begin{code}[language=SQL]
SELECT l.userid
FROM login AS l
WHERE l.sucesso = 0
GROUP BY l.userid
HAVING count(*) > ALL
    (SELECT count(*)
     FROM login AS l1
     WHERE l1.sucesso = 1
         AND l1.userid =l.userid);
		\end{code}
		
Esta query é relativamente simples, assim, a nossa implementação baseou-se em contar o número de sucessos de cada utilizador e retorna-lo se este fosse maior que o número de insucessos do mesmo utilizador.
\end{minipage}

\begin{minipage}{\textwidth}
	\item Quais são os registos que aparecem em todas as paginas de um utilizador?
	
	A query utiliza \textit{ID\_USER} que deve ser substituido pelo id do utilizador desejado (userid):
		\lstset{emph={ID_USER}, emphstyle=\itshape} % da enfase a palavra ID_USER
		\begin{code}[language=SQL]
SELECT r.regcounter
FROM registo AS r
WHERE r.ativo
    AND r.userid = ID_USER
    AND NOT EXISTS
        (SELECT p.pagecounter
         FROM pagina AS p
         WHERE p.userid = ID_USER
             AND r.regcounter NOT IN
                 (SELECT rp.regid
                  FROM reg_pag AS rp
                  WHERE rp.regid = r.regcounter
                      AND rp.pageid = p.pagecounter
                      AND rp.userid = ID_USER
                      AND rp.ativa
                      AND p.ativa
                      AND EXISTS
                          (SELECT tp.typecnt
                           FROM tipo_registo AS tp
                           WHERE tp.typecnt = rp.typeid
                               AND tp.userid = ID_USER
                               AND tp.ativo)))
        		\end{code}
        		
Para esta query seguimos a lógica de que os registos que aparecem em todas as páginas de um utilizador são também os registos para o qual não existe uma página que não tenha esse registo. 

Foi também verificado sempre se as entradas nas tabelas se encontravam activas.
\end{minipage}	

\begin{minipage}{\textwidth}
	\item Quais os utilizadores que têm o maior número médio de registos por página?
	
		\begin{code}[language=SQL]
SELECT rp.userid
FROM reg_pag as rp
WHERE rp.ativa and 
	exists (
		select p.pagecounter 
		from pagina p
		where 	p.userid = rp.userid and
			p.pagecounter = rp.pageid and
			p.ativa) and
	exists (
		select r.regcounter
		from registo r
		where	r.userid = rp.userid and
			r.regcounter = rp.regid and
			r.ativo)
			
GROUP BY rp.userid
HAVING count(*) / count(DISTINCT rp.pageid) >= all
    (SELECT count(*) / count(DISTINCT rp2.pageid)
     FROM reg_pag rp2
     WHERE rp2.ativa and 
		exists (
			select p1.pagecounter 
			from pagina p1
			where 	p1.userid = rp2.userid and
				p1.pagecounter = rp2.pageid and
				p1.ativa) and
		exists (
			select r1.regcounter
			from registo r1
			where	r1.userid = rp2.userid and
				r1.regcounter = rp2.regid and
				r1.ativo)
	     GROUP BY rp2.userid);
		\end{code}
		
Para esta query calculamos a média do número de registos por página de um utilizador que tem a média do número de registos por página maior que a de todos os outros utilizadores. Para tal fomos sempre verificando se as entradas estavam activas.

 
\end{minipage}
\begin{minipage}{\textwidth}
	\item Quais os utilizadores que, em todas as suas páginas, têm registos de todos os tipos de registos que criaram?
	
\begin{code}[language=SQL]
SELECT u.userid,
       u.nome
FROM
  ( SELECT t.userid,
           min(t.num_tipos_pagina) AS minimo
   FROM
     ( SELECT rp.userid,
              rp.pageid,
              count(DISTINCT rp.typeid) AS num_tipos_pagina
      FROM registo r,
           tipo_registo tr,
           pagina p,
           reg_pag rp
      WHERE r.ativo
        AND tr.ativo AND p.ativa AND rp.ativa
        AND r.typecounter=tr.typecnt
        AND r.regcounter=rp.regid
        AND p.pagecounter=rp.pageid
        AND r.userid=tr.userid
        AND p.userid=r.userid
        AND rp.userid=p.userid
      GROUP BY rp.pageid) t
   GROUP BY userid) t,
     utilizador u
WHERE t.userid=u.userid
  AND t.minimo =
    (SELECT count(*)
     FROM tipo_registo tr1
     WHERE tr1.ativo
       AND tr1.userid= t.userid);
\end{code}

Para esta query seguimos a lógica de que os utilizadores que o utilizador tem em todas as suas páginas registos com todos os tipos de registos que criaram se o mínimo de tipos diferentes que tiver em todas as suas páginas for igual ao número de tipo de registos que criou. Para esta query não conseguimos retirar as sub-queries dos FROMs, mesmo após algum esforço e por isso acabamos por ter de deixar como está.

\end{minipage}
\end{enumerate}

\subsection{Notas sobre as Consultas SQL}
\subsubsection{Consulta b -  Quais são os registos que aparecem em todas as páginas de um utilizador?}
Para esta consulta assume-se que se um utilizador não tem páginas então é impossível um registo desse mesmo utilizador aparecer em alguma(s) página(s) (devido à inexistência destas).

Entendeu-se também com esta questão que eram pedidos os registos que aparecem em todas as páginas de um dado utilizador, caso contrário diria "para cada o utilizador". No entanto, como o número de casos em que o facto é verídico é baixo decidimos verificar então quais os registos por todos os utilizadores que verificavam a consulta, e por isso deixamos em baixo a consulta que nos permitiu ver essa informação:

\begin{minipage}{\textwidth}
\begin{code}[language=SQL]
SELECT r_0.userid,
       r_0.regcounter
FROM registo AS r_0
WHERE regcounter IN
        (SELECT r.regcounter
         FROM registo AS r
         WHERE r.ativo
             AND r.userid = r_0.userid
             AND NOT EXISTS
                 ( SELECT p.pagecounter
                  FROM pagina AS p
                  WHERE p.userid = r_0.userid
                      AND r.regcounter NOT IN
                          ( SELECT rp.regid
                           FROM reg_pag AS rp
                           WHERE rp.regid = r.regcounter
                               AND rp.pageid = p.pagecounter
                               AND rp.userid = r_0.userid
                               AND rp.ativa
                               AND p.ativa
                               AND EXISTS
                                   ( SELECT tp.typecnt
                                    FROM tipo_registo AS tp
                                    WHERE tp.typecnt = rp.typeid
                                        AND tp.userid = r_0.userid
                                        AND tp.ativo))))
GROUP BY r_0.userid,
         r_0.regcounter
         
\end{code}

\end{minipage}

\newpage

\section{Restrições de Integridade}
As restrições de integridade foram implementadas com triggers.

Como os triggers acabam por ser todos muito semelhantes apenas apresentaremos aqui um trigger para um insert e um para um update. Os restantes estarão no ficheiro zip da entrega.


\begin{minipage}{\textwidth}
\begin{code}[language=SQL]
drop trigger if exists contador_sequencia_registo_insert;

delimiter $$

    create trigger contador_sequencia_registo_insert before insert on registo
    for each row
    begin
    if (exists(select * from  tipo_registo tr where tr.idseq = new.idseq) or
    	exists(select * from  pagina p where p.idseq = new.idseq) or
    	exists(select * from  campo c where c.idseq = new.idseq) or
    	exists(select * from  registo r where r.idseq = new.idseq) or
    	exists(select * from  valor v where v.idseq = new.idseq))

    then
            call contador_sequencia_registo_insert_trigger();
    end if;
    end$$

delimiter ;

drop trigger if exists contador_sequencia_registo_update;

delimiter $$

    create trigger contador_sequencia_registo_update before update on registo
    for each row
    begin
    if (exists(select * from  tipo_registo tr where tr.idseq = new.idseq) or
    	exists(select * from  pagina p where p.idseq = new.idseq) or
    	exists(select * from  campo c where c.idseq = new.idseq) or
    	exists(select * from  registo r where r.regcounter <> old.regcounter and r.idseq=new.idseq ) or
    	exists(select * from  valor v where v.idseq = new.idseq))

    then
            call contador_sequencia_registo_update_trigger();
    end if;
    end$$

delimiter ;
         
\end{code}

O trigger de insert funciona da seguinte maneira, quando um utilizador tenta inserir uma nova entrada na tabela de registos, se já existir alguma sequencia com o mesmo valor noutra tabela chama uma função inexistente que fará com que o mysql dê erro e aborte a inserção da entrada. Impedindo assim resultados incorrectos

No trigger de update é necessário verificar também se já existe uma sequência nas outras tabelas com a mesma sequência e neste caso deve também verificar se já existe um registo diferente do que se tenta inserir que tenha o mesmo descritivo de sequência do que o que estamos a tentar inserir. Caso aconteça é chamada uma função inexistente para proibir a modificação da tabela.

\end{minipage}
\newpage
\section{Formas Normais}

\begin{itemize}
	\item[(a)]
	A relação utilizador, tem apenas as dependências funcionais (DFs) da forma \textit{X $\to$ A},
	com \textit{A} pertencente aos atributos desta relação, \textit{X  $\subseteq$ \{userid, email\}}
	e \textit{X $\neq \emptyset$} (para além das DFs triviais). Como em todas estas DFs se tem que o determinante é chave,
	a relação utilizador encontra-se na \textit{Boyce-Codd Normal Form} (BCNF).

	\item[(b)]
	Para além das DFs anteriores, a relação tem agora uma nova DF na qual o determinante
	não é chave, mas o dependente é. Assim, a relação utilizador encontra-se na terceira forma
	normal. 
	
	Propomos a seguinte decomposição:
	\begin{flalign*}
		&R_1(\underline{userid}, email, pais, categoria) \\
		&R_2(\underline{userid}, nome, questao1, resposta1, questao2, resposta2).
	\end{flalign*}
	Ambas as relações estão agora na BCNF, pois todas as suas DFs têm chaves como determinantes.
	Para além disso, pelo Teorema de Heath, esta é uma decomposição sem perdas de
	informação, pois \textit{userid} é chave primária nas relações \textit{R}, $R_1$ e $R_2$.
\end{itemize}

\newpage
\section{Índices}
\newpage
\section{Transacções}
Para as transações, de forma a não existirem incoerências devido a acessos simultâneos às mesmas variáveis foi necessário criar transações. As transações fazem o acesso às variáveis dentro de uma transação seja efectuado em regime de exclusividade.

De forma a não proibir o acesso durante muito tempo ao programa, pretendemos ser breves nas transações e apenas realizá-las quando 

Este facto pode ser verificado no código onde uma transação é inicializada com \$connection->beginTransaction() e terminam com \$connection->commit().

\begin{code}[language=PHP]
    <?php

    require "connect.php";

    if (($_SERVER["REQUEST_METHOD"] == "POST")  && ($_POST["nome"] != "") &&
     ($_POST["email"] != "") && ($_POST["questao1"] != "") && 
     ($_POST["questao2"] != "") && ($_POST["resposta1"] != "") 
     && ($_POST["resposta2"] != "")){

        session_start();

        $nome = $_POST["nome"];
        $email = $_POST["email"];
        $password = $_POST["password"];
        $questao1 = $_POST["questao1"];
        $resposta1 = $_POST["resposta1"];
        $questao2 = $_POST["questao2"];
        $resposta2 = $_POST["resposta2"];
        $pais = $_POST["pais"];
        $categoria = $_POST["categoria"];

        (...)
        $connection->beginTransaction();
        
        $query_cria  = (...)
        
        $utilizador_obj = $connection->prepare($query_cria);
    	$utilizador_obj->bindParam(":userid", $userid_aux);
        $utilizador_obj->bindParam(":email", $email_aux);
        $utilizador_obj->bindParam(":nome", $nome_aux);
        $utilizador_obj->bindParam(":password", $password_aux);
        $utilizador_obj->bindParam(":questao1", $questao1_aux);
        $utilizador_obj->bindParam(":resposta1", $resposta1_aux);
        $utilizador_obj->bindParam(":questao2", $questao2_aux);
        $utilizador_obj->bindParam(":resposta2", $resposta2_aux);
        $utilizador_obj->bindParam(":pais", $pais_aux);
        $utilizador_obj->bindParam(":categoria", $categoria_aux);


		(...)

        $connection->commit();
        header("Location: principal.php");
        }

    $connection = null;
    ?>	
\end{code}

\newpage

\section{Data Warehouse}
\newpage

\section{Conclusão}

Concluindo este relatório, achamos que aprendemos imenso com este projecto. Infelizmente tivemos imensos problemas com debugg do PHP e algumas dificuldades com os indices. Pensamos ter ultrapassado a maior parte das nossas dificuldades mas infelizmente não conseguimos resolver todos os problemas do código do projecto.

Iremos tentar termina-lo na mesma, mesmo fora do âmbito da cadeira pois pensamos que nos é de bastante interesse.
\newpage

\end{document}
